\section*{Response to the 4 October 2026 Report on R11}
We thank the referee for distinguishing the policy-specific verification progress from the unresolved methodological comparison. We have revised the original Neural Bellman Operators paper rather than replacing its subject. The response adds direct policy contrasts, a population design, a precise observation implementation, a mechanism-level evaluation bound, independent numerical benchmarks and a single-source reproduction pipeline. The retained applications, earlier adverse results and all historical files remain available. Numerical results below are generated from the fixed final protocol, not entered by hand.

\subsection*{B1 and M4: What does Bellman evaluation contribute?}
The new evaluation--improvement proposition bounds feasible Hamiltonian loss by a quadratic function of critic costate error and approximate actor-improvement error. The accompanying conditional-projection identity separates the variance that regression can remove from the bias it may introduce. We do not claim these classical mathematical principles as new policy-iteration theory, or infer optimizer convergence from them. Instead, we use them to specify direct tests: raw rollout costates without a learned critic, value-only and costate-only regression, one versus five versus ten block updates, and width changes. The primary direct-policy comparison uses the same nonlinear actor without a critic. Its competitive or adverse outcomes are retained. The certificate is correctly described as method-agnostic, while the numerical value added by the Bellman block is measured separately.

\subsection*{B2, M2 and M3: Direct paired contrasts and total work}
The new experiment computes the common-path NBO-minus-direct-policy and NBO-minus-affine statistics themselves. It verifies common initial profiles and innovation hashes and checks that the common reference trajectory cancels. The deterministic bound then contains the two policy transfer errors and conservative statistic-roundoff budgets, without counting the reference-state transfer twice. A clipping-envelope argument and a preallocated finite-family empirical Bernstein calculation provide simultaneous two-sided endpoints. Separate schedule-relative lower endpoints are never subtracted. The comparison now uses equal announced fitting wall budgets, records iteration-completion overshoot and includes validation and checkpointing in the clock. Simulation visits, update counts, verification cost, peak memory and actor deployment timing remain separate. Full-path checkpoint frontiers are retained rather than relying on low-path diagnostic endpoints. A method advantage is stated only to the extent supported by its direct endpoint; a favorable sample mean alone is insufficient.

\subsection*{B3: Accuracy relative to the economic optimum}
We retain the actual capital regret endpoints rather than relabeling safe improvement as a near-optimal HJB solve. We add an independent monotone implicit finite-difference reference for the nonlinear one-state member of the same capital family. It solves the full primitive action-box problem without a neural critic, evaluates fixed learned policies, and reports nested spatial/time grids, a wider domain, equation residuals and policy-iteration convergence. These finite-grid comparisons provide an absolute numerical policy-loss diagnostic. They do not by themselves enclose the continuous multidimensional optimum, and the text makes that distinction explicit. The high-dimensional continuous guarantee remains the verified policy-specific bound rather than an extrapolation from the scalar computation.

\subsection*{B4 and M8: A predefined population and visible failures}
The final protocol defines nine initial balance-sheet profiles through aggregate log-capital and dispersion. Initial profiles are drawn independently with fixed equal probabilities on final paths. The theorem now gives integrated performance for this population by averaging the deterministic transfer accounts and using global profile-valid clipping bounds. It does not claim uniform control over a continuum of initial states. Primary population and origin comparisons are both reported. Severe mean and dispersion stresses appear in the main article for both nonlinear methods, including negative and inconclusive outcomes. The population is a transparent model-based design, not an asserted empirical distribution.

\subsection*{B5: Observation model and componentwise clipping}
We now prove implementation from continuously observed noiseless capital histories, known coefficients, own actions and independent private randomization. Subtracting the integrated known drift from the state path gives the aggregate innovation. Because the original implementation clips $d+1$ drivers separately while the state has dimension $d$, a right inverse alone would not preserve its law. We add the missing null-space component using one independent scalar Brownian randomization. The reconstructed vector has identity bracket, has the correct aggregate state increment and gives the same closed-loop law, including clipping. This is a distributional implementation theorem, not pathwise recovery of the original latent null driver. It does not solve discrete noisy observation, and the exact drift-integral observation premise is explicit. Numerical covariance, right-inverse and null-space regression checks accompany the proof.

\subsection*{B6 and M10: Economic content, scope and title}
The title remains Neural Bellman Operators and the original economic program is retained. The new main-text organization separates the evaluation map, the conditional action-loss relation, the policy-verification theorem and evidence of comparative performance. We do not attribute a generic Bellman principle or universal neural advantage to the verification result. To make the economic scale operational, we add a self-financed management-fee experiment: a fixed fraction of gross withdrawals pays for administration, while net consumption enters logarithmic felicity. The exact fee identity transforms verified utility differences into explicit supported or unsupported adoption thresholds. This is a stated quantitative decision in the stylized economy, not a calibrated welfare claim. Earlier models and evidence remain in the supplement instead of being deleted to obtain a narrower-looking paper.

\subsection*{M1: Raw averages and continuous-time endpoints}
All new tables use the label mean paired statistic for the unadjusted average. The continuous-time lower and upper endpoints separately include diffusion transfer, empirical Bernstein uncertainty, clipping tails and arithmetic. Retained R11 tables are explicitly identified as historical and their mean labels receive this same interpretation. A confidence endpoint is never replaced by the raw average.

\subsection*{M5: Training randomness and selection}
The final study uses ten new fixed initialization/minibatch streams. Training and validation are independent of the final noise bank, and selection uses validation only. The final protocol fixes all dimensions, profiles, radii, contrasts and auxiliary variants before final fitting. Seed-level outcomes and work counts are retained. Their variation is descriptive of the recorded streams; we do not claim a sampling theorem for an unspecified population of hyperparameters or future optimizers. Development and final noise banks differ, and the timing-based stopping rule is documented as hardware-dependent.

\subsection*{M6: One clean final source and execution}
The publication workflow checks out one fixed source commit for every numerical worker and for audit and compilation. Every primary and auxiliary fit, population evaluation, stress, direct contrast and reference calculation is produced in that run. The final numerical stage contains no recovery layer and does not import selected outcomes from a previous run. Development failures are documented separately. The evidence commit contains the generated records and papers, the source hash, workflow identity, environment data and hash manifests; no source-manifest exception is used to change a fitted or evaluated object after testing. Earlier source-layout tests retain their historical context, while a separate full-tree preservation audit checks this revision against the reviewed snapshot.

\subsection*{M7: Fresh radius fits}
Both nonlinear methods are refitted at radii 0.05 and 0.15 under the announced fitting budget. These are distinct trained objects with their own validation selection, weights, cost and final certificates, not scaled versions of the radius-0.10 network. Their results and the other ablations are retained seed by seed.

\subsection*{M9: Repository entry point}
The root README now identifies R13, its review target, authoritative manuscripts, source/evidence workflow, reading order and current reproduction commands. The former README is archived exactly. Generated audit counts and source identities are linked rather than left at the obsolete September 28 numbers.

\section*{Executed evidence}
\input{revisions/2026-10-04-r13/manuscript/results_summary.tex}
The accompanying audit records software completion separately from positive improvement, method superiority, fee support and absolute policy accuracy. All statements remain subject to their declared economic, observation, numerical and statistical scope; repository execution is not a journal decision.
